\documentclass[11pt]{article}
\usepackage[utf8]{inputenc}
\usepackage[T1]{fontenc}
\usepackage{amsmath,amssymb,mathtools}
\usepackage[margin=1in]{geometry}
\usepackage{hyperref}
\usepackage{array}
\usepackage{parskip}
\setlength{\parindent}{0pt}
\hypersetup{colorlinks=true,linkcolor=blue,urlcolor=blue}
\title{Lecture 8 — Differential / Difference Equations and Singularity Functions}
\author{}
\date{}
\begin{document}
\maketitle

\textbf{Course:} CEC 315 — Signals and Systems (Spring 2026)
\textbf{Instructor:} Rogelio Gracia Otalvaro
\textbf{Source PDF:} \texttt{all\_lectures/cec315-lctr08-diff-eqns-singularity.pdf}
\textbf{Exam coverage:} Exam 1

\textbf{Lctr 8: Differential/Difference Equations and Singularity Functions}

Rogelio Gracia Otalvaro

Spring 2026

\textbf{Focus:} Sections 2.4 and 2.5 (Pages 116 to 140)

\section*{8.1 Introduction: Why Differential and Difference Equations?}

\begin{quote}
\textbf{Why This Matters}

So far, we've characterized LTI systems by their \textbf{impulse response} $h(t)$ or $h[n]$. But how do we know what $h(t)$ is in the first place? For many physical systems (circuits, mechanical systems, etc.), the system is \textbf{naturally} described by a \textbf{differential equation} (CT) or \textbf{difference equation} (DT). This lecture shows how to connect these descriptions and find the impulse response!
\end{quote}

\textbf{Two ways to describe an LTI system:}

\begin{center}
\begin{tabular}{|l|l|}
\hline
\textbf{Impulse Response} & \textbf{Differential / Difference Equation} \\
\hline
$y(t) = x(t)*h(t)$ & $\dfrac{dy}{dt} + a\,y = b\,x$ \\
\hline
Convolution integral & Derivative relationship \\
\hline
\end{tabular}
\end{center}

Both describe the \textbf{same} system! We'll learn to go from one to the other.

\section*{8.2 First-Order Differential Equations (Continuous-Time)}

\subsection*{8.2.1 The Standard Form}

The simplest and most common CT LTI system is described by a \textbf{first-order} equation:

\[\boxed{\;\frac{dy(t)}{dt} + a\,y(t) = b\,x(t)\;}\]

where:

\begin{itemize}
  \item $y(t)$ = output signal
  \item $x(t)$ = input signal
  \item $a,\,b$ = constants (system parameters)
\end{itemize}

\textbf{Physical example.} An RC circuit has the equation:

\[\frac{d v_C(t)}{d t} + \frac{1}{RC}\,v_C(t) = \frac{1}{RC}\,v_{\text{in}}(t)\]

Here $a = 1/(RC)$ is the "time constant" (inverse of) of the circuit.

\subsection*{8.2.2 Solving First-Order Equations: Step by Step}

To solve $\tfrac{dy}{dt} + a\,y = b\,x(t)$, we use the \textbf{integrating factor method} (or, equivalently, split the solution into homogeneous + particular parts).

\textbf{Step 1:} Identify $a$ and $b$ from the equation.

\textbf{Step 2:} The \textbf{homogeneous solution} (when $x(t) = 0$) is:

\[y_h(t) = C\,e^{-at}\]

where $C$ is determined by initial conditions.

\textbf{Step 3:} Find the \textbf{particular solution} $y_p(t)$ based on the form of $x(t)$ (try a form similar to the input).

\textbf{Step 4:} The \textbf{total solution} is:

\[y(t) = y_h(t) + y_p(t) = C\,e^{-at} + y_p(t)\]

\textbf{Step 5:} Apply initial conditions to find $C$.

\begin{quote}
\textbf{Key Insight — Initial Rest}

For a \textbf{causal LTI system at initial rest}, the initial condition is $y(0^-) = 0$. This means the system has no stored energy before the input is applied. "Initial rest" is the convention we use throughout this course.
\end{quote}

\subsection*{8.2.3 Worked Example 1: Finding the Impulse Response}

\textbf{Problem:} Find the impulse response $h(t)$ of the system

\[\frac{dy(t)}{dt} + 2\,y(t) = x(t).\]

\textbf{Solution.} Set $x(t) = \delta(t)$ (unit impulse) and solve for $y(t) = h(t)$.

\textbf{Step 1: Identify parameters.} $a = 2$, $b = 1$.

\textbf{Step 2: For $t>0$,} the impulse has already passed and the equation becomes homogeneous:

\[\frac{dh(t)}{dt} + 2\,h(t) = 0\qquad t > 0\]

with solution $h(t) = C\,e^{-2t}$ for $t > 0$.

\textbf{Step 3: Find $C$ using the impulse's effect.} The impulse $\delta(t)$ causes an instantaneous jump in the output. Integrate the original equation across $t=0$:

\[\int_{0^-}^{0^+} \frac{dh}{dt}\,dt + 2\int_{0^-}^{0^+} h(t)\,dt = \int_{0^-}^{0^+} \delta(t)\,dt\]

The first integral is $h(0^+) - h(0^-)$. The second integral is zero (because $h$ is finite over an infinitesimal interval). The right side is 1. So:

\[h(0^+) - h(0^-) = 1\]

With initial rest ($h(0^-) = 0$) we get $h(0^+) = 1$. From $h(t) = C\,e^{-2t}$ at $t = 0^+$: $1 = C\,e^{0} = C$, so $C = 1$.

\textbf{Step 4: Write the complete answer.}

\[\boxed{\;h(t) = e^{-2t}\,u(t)\;}\]

The $u(t)$ ensures the response is zero for $t<0$ (causality).

\begin{quote}
\textbf{Key Insight — Quick Formula}

For the first-order system $\tfrac{dy}{dt} + a\,y = b\,x$, the impulse response is:

\[h(t) = b\,e^{-at}\,u(t)\]

This is a decaying exponential if $a > 0$ (stable system).
\end{quote}

\subsection*{8.2.4 Worked Example 2: Response to a Step Input}

\textbf{Problem:} Find the output $y(t)$ of the system $\tfrac{dy}{dt} + 2y = 3x$ when $x(t) = u(t)$, with $y(0^-) = 0$.

\textbf{Step 1:} $a = 2$, $b = 3$.

\textbf{Step 2:} For $t \geq 0$, the equation becomes:

\[\frac{dy}{dt} + 2y = 3\cdot 1 = 3\]

\textbf{Step 3:} Homogeneous solution: $y_h(t) = C\,e^{-2t}$.

\textbf{Step 4:} Particular solution. Since the RHS is a constant (3), try $y_p = K$ (constant):

\[\frac{d K}{dt} + 2K = 3\;\Rightarrow\;0 + 2K = 3\;\Rightarrow\;K = \frac{3}{2}\]

\textbf{Step 5:} Total solution:

\[y(t) = C\,e^{-2t} + \frac{3}{2}\quad\text{for } t \geq 0\]

\textbf{Step 6:} Apply initial condition. At $t=0$: $y(0) = 0$ (initial rest, since there is no impulse at $t=0$):

\[0 = C\,e^{0} + \frac{3}{2} = C + \frac{3}{2}\;\Rightarrow\;C = -\frac{3}{2}\]

\textbf{Final answer:}

\[\boxed{\;y(t) = \frac{3}{2}\left(1 - e^{-2t}\right)u(t)\;}\]

\textbf{Interpretation:} the output starts at 0 and exponentially approaches the steady-state value $3/2$. Note that this matches the step response $s(t) = h(t)*u(t)$ you would obtain from the impulse response $h(t) = 3\,e^{-2t}\,u(t)$.

\section*{8.3 First-Order Difference Equations (Discrete-Time)}

\subsection*{8.3.1 The Standard Form}

The discrete-time equivalent is a \textbf{first-order difference equation}:

\[\boxed{\;y[n] - \alpha\,y[n-1] = b\,x[n]\;}\]

or equivalently (rearranging):

\[y[n] = \alpha\,y[n-1] + b\,x[n]\]

This is a \textbf{recursive} equation: the current output depends on the previous output.

\begin{quote}
\textbf{Key Insight — Recursive vs. Non-recursive}

- \textbf{Recursive (IIR):} output depends on past outputs. Has an \textbf{I}nfinite \textbf{I}mpulse \textbf{R}esponse.
- \textbf{Non-recursive (FIR):} output depends only on inputs. Has a \textbf{F}inite \textbf{I}mpulse \textbf{R}esponse.
\end{quote}

\subsection*{8.3.2 Solving by Iteration}

The easiest way to solve difference equations is by \textbf{iteration}: compute values one by one, starting from initial rest.

\textbf{Worked Example 3 — Finding Impulse Response by Iteration.}

\textbf{Problem:} Find the impulse response $h[n]$ of $y[n] = \tfrac{1}{2}\,y[n-1] + x[n]$.

\textbf{Solution.} Set $x[n] = \delta[n]$ and compute $h[n]$ for each $n$, starting from initial rest ($y[n] = 0$ for $n < 0$).

\[n = 0:\quad h[0] = \tfrac{1}{2}\,h[-1] + \delta[0] = \tfrac{1}{2}(0) + 1 = 1\]
\[n = 1:\quad h[1] = \tfrac{1}{2}\,h[0] + \delta[1] = \tfrac{1}{2}(1) + 0 = \tfrac{1}{2}\]
\[n = 2:\quad h[2] = \tfrac{1}{2}\,h[1] + \delta[2] = \tfrac{1}{2}\!\left(\tfrac{1}{2}\right) + 0 = \tfrac{1}{4}\]
\[n = 3:\quad h[3] = \tfrac{1}{2}\,h[2] + \delta[3] = \tfrac{1}{2}\!\left(\tfrac{1}{4}\right) + 0 = \tfrac{1}{8}\]

\textbf{Pattern:} $h[n] = (1/2)^n$ for $n \geq 0$.

\[\boxed{\;h[n] = \left(\tfrac{1}{2}\right)^n u[n]\;}\]

\begin{quote}
\textbf{Key Insight — Quick Formula}

For the first-order system $y[n] = \alpha\,y[n-1] + b\,x[n]$, the impulse response is:

\[h[n] = b\,\alpha^n\,u[n]\]

The system is \textbf{stable} if $|\alpha| < 1$ (decaying geometric sequence).
\end{quote}

\subsection*{8.3.3 Worked Example 4: Step Response by Iteration}

\textbf{Problem:} Find the step response $s[n]$ of $y[n] = \tfrac{1}{2}y[n-1] + x[n]$ when $x[n] = u[n]$.

\textbf{Solution (by iteration).}

\[n = 0:\quad s[0] = \tfrac{1}{2}s[-1] + u[0] = 0 + 1 = 1\]
\[n = 1:\quad s[1] = \tfrac{1}{2}s[0] + u[1] = \tfrac{1}{2}(1) + 1 = \tfrac{3}{2}\]
\[n = 2:\quad s[2] = \tfrac{1}{2}s[1] + u[2] = \tfrac{1}{2}\!\left(\tfrac{3}{2}\right) + 1 = \tfrac{7}{4}\]
\[n = 3:\quad s[3] = \tfrac{1}{2}s[2] + u[3] = \tfrac{1}{2}\!\left(\tfrac{7}{4}\right) + 1 = \tfrac{15}{8}\]

\textbf{Pattern:}

\[s[n] = 2 - \left(\tfrac{1}{2}\right)^n = 2\!\left(1 - \left(\tfrac{1}{2}\right)^{n+1}\right)\]

The output approaches 2 as $n\to\infty$ (steady state).

\section*{8.4 Stability From Differential / Difference Equations}

\begin{quote}
\textbf{Key Insight — Stability Check}

- \textbf{CT:} $\tfrac{dy}{dt} + a\,y = b\,x$ is \textbf{stable} if $a > 0$ (exponential decays).
- \textbf{DT:} $y[n] = \alpha\,y[n-1] + b\,x[n]$ is \textbf{stable} if $|\alpha| < 1$ (geometric series converges).
\end{quote}

\textbf{Examples.}

\begin{itemize}
  \item $\frac{dy}{dt} + 3y = x\;\Rightarrow\;a = 3 > 0\;\Rightarrow$ \textbf{stable}.
  \item $\frac{dy}{dt} - 2y = x\;\Rightarrow\;a = -2 < 0\;\Rightarrow$ \textbf{unstable} (grows exponentially!).
  \item $y[n] = 0.5\,y[n-1] + x[n]\;\Rightarrow\;|\alpha| = 0.5 < 1\;\Rightarrow$ \textbf{stable}.
  \item $y[n] = 2\,y[n-1] + x[n]\;\Rightarrow\;|\alpha| = 2 > 1\;\Rightarrow$ \textbf{unstable}.
\end{itemize}

\section*{8.5 Singularity Functions}

\begin{quote}
\textbf{Why This Matters}

Singularity functions ($\delta(t)$, $u(t)$, and their derivatives / integrals) form a \textbf{family} of related functions. Understanding their relationships helps us analyze systems and solve differential equations more easily.
\end{quote}

\subsection*{8.5.1 The Family of Singularity Functions}

The singularity family is connected by differentiation (moving right) and integration (moving left):

\[r(t) = t\,u(t) \;\xrightarrow{d/dt}\; u(t) \;\xrightarrow{d/dt}\; \delta(t) \;\xrightarrow{d/dt}\; \delta'(t)\]

Going the other way (by integration):

\[\text{Ramp } r(t) \;\xleftarrow{\int}\; \text{Step } u(t) \;\xleftarrow{\int}\; \text{Impulse } \delta(t) \;\xleftarrow{\int}\; \text{Doublet } \delta'(t)\]

\textit{[Figure: The four members of the singularity family displayed as a horizontal chain. From left to right: ramp $r(t) = t\,u(t)$, step $u(t)$, impulse $\delta(t)$, doublet $\delta'(t)$. Arrows pointing right are labeled $d/dt$; arrows pointing left are labeled $\int$.]}

\subsection*{8.5.2 Key Relationships}

\textbf{Derivative relationships:}

\[\frac{d}{dt}[r(t)] = u(t),\qquad \frac{d}{dt}[u(t)] = \delta(t),\qquad \frac{d}{dt}[\delta(t)] = \delta'(t)\]

\textbf{Integral relationships:}

\[\int_{-\infty}^{t}\delta(\tau)\,d\tau = u(t),\qquad \int_{-\infty}^{t}u(\tau)\,d\tau = r(t) = t\,u(t)\]

\begin{quote}
\textbf{Key Insight}

Remember: moving \textbf{right} in the diagram = taking derivatives. Moving \textbf{left} = integrating.
\end{quote}

\subsection*{8.5.3 The Ramp Function}

The \textbf{ramp function} is defined as:

\[r(t) = t\,u(t) = \begin{cases} 0, & t < 0 \\ t, & t \geq 0 \end{cases}\]

It is the integral of the unit step (accumulation of a constant).

\subsection*{8.5.4 The Unit Doublet}

The \textbf{unit doublet} $\delta'(t)$ (also written $u_1(t)$) is the derivative of the impulse:

\[\delta'(t) = \frac{d\delta(t)}{dt}\]

\textbf{Property — Convolving with the doublet takes the derivative:}

\[x(t)*\delta'(t) = \frac{d x(t)}{d t}\]

\section*{8.6 Useful Properties of Singularity Functions}

\subsection*{8.6.1 Scaling the Impulse}

\[\boxed{\;\delta(at) = \frac{1}{|a|}\,\delta(t)\;}\]

\textbf{Example.} $\delta(2t) = \tfrac{1}{2}\,\delta(t)$.

\textbf{Why?} The impulse is "compressed" by factor $a$, but its area must remain 1, so its "height" increases by $|a|$ — hence the $1/|a|$ scale.

\subsection*{8.6.2 Impulse Times a Function}

\[\boxed{\;x(t)\,\delta(t - t_0) = x(t_0)\,\delta(t - t_0)\;}\]

The impulse "samples" the value of $x(t)$ at $t = t_0$.

\textbf{Example.} $(t^2 + 3)\,\delta(t - 2) = (4 + 3)\,\delta(t - 2) = 7\,\delta(t - 2)$.

\begin{quote}
\textbf{Common Mistake}

Don't confuse this with the \textbf{sifting property}! This is multiplication, not integration:

- Multiplication: $x(t)\,\delta(t-t_0) = x(t_0)\,\delta(t-t_0)$ (result is a \textbf{scaled impulse}).
- Sifting (integration): $\int x(t)\,\delta(t-t_0)\,dt = x(t_0)$ (result is a \textbf{number}).
\end{quote}

\section*{8.7 Worked Example 5: Using Singularity Functions}

\textbf{Problem:} Find the derivative of $x(t) = (t + 1)\,u(t)$.

\textbf{Solution.} Use the product rule:

\[\frac{d}{dt}\!\left[(t+1)\,u(t)\right] = \frac{d(t+1)}{dt}\cdot u(t) + (t+1)\cdot \frac{d u(t)}{dt}\]

We know $du(t)/dt = \delta(t)$, so:

\[= 1\cdot u(t) + (t+1)\cdot\delta(t)\]

Using the multiplication rule $(t+1)\,\delta(t) = (0+1)\,\delta(t) = \delta(t)$:

\[\boxed{\;\frac{d}{dt}\!\left[(t+1)\,u(t)\right] = u(t) + \delta(t)\;}\]

\textbf{Interpretation.} The derivative has:

\begin{itemize}
  \item A \textbf{step} $u(t)$: the slope of the ramp part (the function grows with slope 1 for $t \geq 0$).
  \item An \textbf{impulse} $\delta(t)$: the "jump" at $t=0$ from $0$ up to $1$ (the function jumps from 0 to 1 at the origin because $(0+1)\cdot u(0^+) = 1$).
\end{itemize}

\section*{8.8 Summary and Key Formulas}

\begin{center}
\begin{tabular}{|l|l|}
\hline
\textbf{Concept} & \textbf{Formula} \\
\hline
First-order CT system & $\dfrac{dy}{dt} + a\,y = b\,x$ \\
\hline
CT impulse response & $h(t) = b\,e^{-at}\,u(t)$ \\
\hline
CT stability condition & $a > 0$ \\
\hline
First-order DT system & $y[n] = \alpha\,y[n-1] + b\,x[n]$ \\
\hline
DT impulse response & $h[n] = b\,\alpha^n\,u[n]$ \\
\hline
DT stability condition & $|\alpha| < 1$ \\
\hline
Derivative of step & $\dfrac{d u(t)}{dt} = \delta(t)$ \\
\hline
Integral of impulse & $\displaystyle\int_{-\infty}^{t}\delta(\tau)\,d\tau = u(t)$ \\
\hline
Impulse scaling & $\delta(at) = \dfrac{1}{|a|}\,\delta(t)$ \\
\hline
Impulse sampling & $x(t)\,\delta(t - t_0) = x(t_0)\,\delta(t - t_0)$ \\
\hline
Ramp function & $r(t) = t\,u(t)$ \\
\hline
Doublet derivative & $x(t)*\delta'(t) = dx(t)/dt$ \\
\hline
\end{tabular}
\end{center}

\section*{8.9 Common Mistakes to Avoid}

\begin{enumerate}
  \item \textbf{Forgetting the $u(t)$ or $u[n]$:} impulse responses of causal systems must include the unit step to ensure $h(t) = 0$ for $t<0$ (or $h[n] = 0$ for $n<0$).
  \item \textbf{Sign errors in the exponent:} for $\tfrac{dy}{dt} + a\,y = b\,x$, the impulse response has $e^{-at}$, \textbf{not} $e^{+at}$. Watch the signs!
  \item \textbf{Confusing stability conditions:}
\end{enumerate}
\begin{itemize}
  \item \textbf{CT:} need $a > 0$ (positive coefficient of $y$).
  \item \textbf{DT:} need $|\alpha| < 1$ (magnitude less than 1).
\end{itemize}
\begin{enumerate}
  \item \textbf{Iteration errors in DT:} always start from initial rest ($y[n] = 0$ for $n < 0$) and compute values in order.
  \item \textbf{Derivative of $u(t)$:} remember $du(t)/dt = \delta(t)$, \textbf{not} $u(t)$ or 0!
\end{enumerate}

\begin{quote}
\textbf{Key Insight — Study Tip}

Practice the iteration method for DT systems — it builds intuition for how recursive systems work and is a reliable way to check your closed-form answers!
\end{quote}

Rogelio Gracia Otalvaro

\section*{Practice Problems Summary}

\begin{enumerate}
  \item \textbf{Worked Example 1 — CT impulse response:} $\frac{dy}{dt} + 2y = x$ has impulse response $h(t) = e^{-2t}u(t)$, obtained by solving the homogeneous equation for $t>0$ and applying the jump condition $h(0^+) - h(0^-) = 1$.
  \item \textbf{Worked Example 2 — CT step response:} $\frac{dy}{dt} + 2y = 3x$ with $x(t) = u(t)$ gives $y(t) = \tfrac{3}{2}(1 - e^{-2t})\,u(t)$; steady state $= 3/2$.
  \item \textbf{Worked Example 3 — DT impulse response by iteration:} $y[n] = \tfrac{1}{2}y[n-1] + x[n]$ yields $h[n] = (1/2)^n u[n]$ (iterate: 1, 1/2, 1/4, 1/8, …).
  \item \textbf{Worked Example 4 — DT step response by iteration:} Same system with $x[n] = u[n]$ gives $s[n] = 2(1 - (1/2)^{n+1})$; iterates: 1, 3/2, 7/4, 15/8, … approaching 2.
  \item \textbf{Stability examples:} $\frac{dy}{dt} + 3y = x$ stable; $\frac{dy}{dt} - 2y = x$ unstable; $y[n] = 0.5 y[n-1] + x[n]$ stable; $y[n] = 2 y[n-1] + x[n]$ unstable.
  \item \textbf{Singularity family relationships:} $\frac{d}{dt}[r(t)] = u(t)$, $\frac{d}{dt}[u(t)] = \delta(t)$, $\frac{d}{dt}[\delta(t)] = \delta'(t)$; integrals invert these.
  \item \textbf{Impulse scaling:} $\delta(2t) = \tfrac{1}{2}\delta(t)$.
  \item \textbf{Impulse times a function:} $(t^2 + 3)\delta(t - 2) = 7\,\delta(t - 2)$.
  \item \textbf{Worked Example 5 — Derivative of $(t+1)u(t)$:} Product rule gives $u(t) + (t+1)\delta(t) = u(t) + \delta(t)$; the step captures the slope of the ramp and the impulse captures the jump at $t=0$.
\end{enumerate}

\end{document}
